\section{The estimates still needed}
\label{sec:remaining-estimates}

The analytic reductions identify concrete sufficient conditions for GRH.
One is square-root cancellation, with an arbitrarily small power loss,
in the smoothed M\"obius sums of
Theorem~\ref{qr-thm-mellin}, for every fixed primitive character and
every fixed smooth test. The critical-strip reduction then applies.
This is a sufficient condition proved here, not an estimate established
here.

A second route uses the weighted prime sums of the preceding sections.
The announced source constructs upper bounds on long intervals
$[Y,Y^{\lambda}]$ with $\lambda>1$, for suitable finite families of
characters. Three matters must remain explicit: which characters
belong to the family, whether the family is fixed as $Y$ grows, and
whether constants and exponents are uniform in the parameters that
will later vary. A bound only on such intervals cannot be replaced
silently by an eventual bound for every real $x$.

The current source audit locates the input to these bounds in theta
lifts, integral structures, and quantitative bounds for arithmetic
cohomology groups. Their definitions, compatibility conditions, and
parameter ranges have not yet been reduced to a complete checked
ordinary-mathematics proof in this manuscript. No result in the paper
removes this obligation.

The second-moment calculation gives a useful restriction on possible
simplifications. The bound of Proposition~\ref{qr-prop-power-rows},
with its present number of selected indices and its present error,
cannot reach exponent $1/2$ by choosing the length differently.
The best formal exponent for sixth powers is $6/7$ if all positive
length exponents are hypothetically allowed; in the range actually
stated in the October source the limiting exponent is $11/12$.
To reach GRH by that route, an improved arithmetic estimate or a
different use of the family is necessary.

Similarly, the finite descent and the uniform-continuation argument
are useful only with the quantifiers proved in their hypotheses.
A fixed positive gap cannot be replaced by a gap tending to zero
without bounds on the resulting constants and number of steps.
Separate positive improvements for infinitely many characters do not
produce a common improvement. These limitations specify mathematical
work to be done; they are not counterexamples to GRH or to the
announced proof.
